\section{MAVLink-Based Position Control in Guided Mode}

\subsection{Overview}

MAVLink (Micro Air Vehicle Link) is a lightweight communication protocol used for interaction between a ground control system and an autopilot. In guided flight modes (e.g., \texttt{GUIDED} in ArduPilot), MAVLink enables external programs to command the vehicle’s motion in real time.

A common approach for precise control, especially in indoor or GPS-denied environments, is to use \textbf{local position control} rather than global GPS coordinates. This is achieved using the \texttt{LOCAL\_NED} reference frame and corresponding MAVLink messages.

\subsection{LOCAL\_NED Coordinate Frame}

The \texttt{MAV\_FRAME\_LOCAL\_NED} frame is a Cartesian coordinate system defined as:

\begin{itemize}
    \item $x$ axis: North
    \item $y$ axis: East
    \item $z$ axis: Down
\end{itemize}

A key characteristic of this frame is that altitude is represented as a \textbf{negative value in the $z$ direction}. Therefore, commanding a position above the origin requires:

\[
z = -h
\]

where $h$ is the desired altitude above the reference point.

The origin of the local frame is typically established at the vehicle’s position when the system initializes or arms.

\subsection{Obtaining Local Position}

The vehicle’s state in the local frame is provided by the MAVLink message:

\begin{verbatim}
LOCAL_POSITION_NED
\end{verbatim}

This message includes:
\begin{itemize}
    \item Position: $(x, y, z)$ in meters
    \item Velocity: $(v_x, v_y, v_z)$ in m/s
\end{itemize}

This information is essential for:
\begin{itemize}
    \item Feedback control
    \item Monitoring vehicle motion
    \item Generating relative trajectories
\end{itemize}

\subsection{Flight Mode and Arming}

Before sending position commands, the vehicle must be placed into a guided control mode:

\begin{verbatim}
set_mode("GUIDED")
\end{verbatim}

This enables the autopilot to accept external setpoint commands.

\subsubsection{Arming}

The motors are armed using the MAVLink command:

\begin{verbatim}
MAV_CMD_COMPONENT_ARM_DISARM
\end{verbatim}

This is typically sent using a \texttt{COMMAND\_LONG} message, with parameter 1 set to 1 (arm) or 0 (disarm). Confirmation is usually obtained via heartbeat messages.

\subsection{Takeoff Procedure}

Takeoff is commanded using:

\begin{verbatim}
MAV_CMD_NAV_TAKEOFF
\end{verbatim}

The desired altitude is provided as a parameter. The vehicle then ascends to this altitude under autopilot control. During this phase, altitude can be monitored using:

\begin{verbatim}
GLOBAL_POSITION_INT
\end{verbatim}

which provides relative altitude above the home position.

\subsection{Position Control Using SET\_POSITION\_TARGET\_LOCAL\_NED}

The primary MAVLink message for position control is:

\begin{verbatim}
SET_POSITION_TARGET_LOCAL_NED
\end{verbatim}

This message allows control over:
\begin{itemize}
    \item Position $(x, y, z)$
    \item Velocity $(v_x, v_y, v_z)$
    \item Acceleration $(a_x, a_y, a_z)$
    \item Yaw and yaw rate
\end{itemize}

\subsubsection{Type Mask Logic}

A \textbf{type mask} is used to specify which fields in the message should be ignored. Each bit in the mask corresponds to a specific field.

For position-only control, the mask is configured to ignore velocity, acceleration, and yaw rate:

\begin{verbatim}
type_mask = (
    VX_IGNORE | VY_IGNORE | VZ_IGNORE |
    AX_IGNORE | AY_IGNORE | AZ_IGNORE |
    YAW_RATE_IGNORE
)
\end{verbatim}

If yaw is not controlled:

\begin{verbatim}
type_mask |= YAW_IGNORE
\end{verbatim}

\paragraph{Interpretation:}

\begin{itemize}
    \item Position is controlled directly
    \item Velocity and acceleration are handled internally by the autopilot
    \item Yaw can be either fixed or ignored
\end{itemize}

This simplifies control and leverages the autopilot’s internal controllers for smooth motion.

\subsubsection{Sending Position Setpoints}

The command is sent using:

\begin{verbatim}
set_position_target_local_ned_send(...)
\end{verbatim}

with:
\begin{itemize}
    \item Desired position $(x, y, z)$
    \item Type mask defining active fields
    \item Optional yaw angle
\end{itemize}

\subsection{Continuous Setpoint Streaming}

In guided mode, position commands must be sent at a fixed rate (typically 5–20 Hz). This is necessary because:

\begin{itemize}
    \item The autopilot expects continuous updates
    \item A lack of commands may cause the vehicle to hold position or trigger failsafe behavior
\end{itemize}

Thus, even when holding a fixed position, the same setpoint is repeatedly transmitted.

\subsection{Landing Procedure}

Landing is typically performed by switching to an autonomous landing mode:

\begin{verbatim}
set_mode("LAND")
\end{verbatim}

In this mode, the autopilot handles descent, velocity control, and motor shutdown.

\subsection{Code Examples and Practical Execution}

This subsection shows a minimal workflow for issuing MAVLink position commands from Python using \texttt{pymavlink}. 

\subsubsection{Connecting to the MAVLink stream}

A MAVLink client first needs to connect to a UDP endpoint:

\begin{verbatim}
from pymavlink import mavutil

master = mavutil.mavlink_connection('udpin:127.0.0.1:14551')
master.wait_heartbeat()

print("Connected")
print(f"Heartbeat from system {master.target_system} "
      f"component {master.target_component}")
\end{verbatim}

The call to \texttt{wait\_heartbeat()} is important because it confirms that communication with the autopilot has been established and also provides the target system and component identifiers used in later commands.

\subsubsection{Flight mode, arming, and takeoff}

Before sending local position instructions, the vehicle must be placed in a guided mode, armed, and lifted to a safe altitude. The following example reflects the same sequence.

\begin{verbatim}
def set_mode(master, mode_name):
    mode_map = master.mode_mapping()
    if mode_map is None or mode_name not in mode_map:
        raise RuntimeError(f"Mode {mode_name} not available")

    mode_id = mode_map[mode_name]
    master.mav.set_mode_send(
        master.target_system,
        mavutil.mavlink.MAV_MODE_FLAG_CUSTOM_MODE_ENABLED,
        mode_id
    )

def arm_vehicle(master):
    master.mav.command_long_send(
        master.target_system,
        master.target_component,
        mavutil.mavlink.MAV_CMD_COMPONENT_ARM_DISARM,
        0,
        1, 0, 0, 0, 0, 0, 0
    )

def takeoff(master, target_alt_m):
    master.mav.command_long_send(
        master.target_system,
        master.target_component,
        mavutil.mavlink.MAV_CMD_NAV_TAKEOFF,
        0,
        0, 0, 0, 0,
        0, 0,
        target_alt_m
    )
\end{verbatim}

A typical usage sequence is then:

\begin{verbatim}
set_mode(master, "GUIDED")
arm_vehicle(master)
takeoff(master, 1.5)
\end{verbatim}

The takeoff altitude is then checked using \texttt{GLOBAL\_POSITION\_INT}, and the code waits until the vehicle reaches approximately \(95\%\) of the requested height before continuing. 

\subsubsection{Reading the current local position}

After takeoff, the local position can be obtained from \texttt{LOCAL\_POSITION\_NED}:

\begin{verbatim}
def get_current_local_position(master, timeout=5):
    import time
    deadline = time.time() + timeout

    while time.time() < deadline:
        msg = master.recv_match(
            type=['LOCAL_POSITION_NED', 'STATUSTEXT'],
            blocking=True,
            timeout=1
        )

        if not msg:
            continue

        if msg.get_type() == 'STATUSTEXT':
            print("STATUSTEXT:", msg.text)
            continue

        return msg.x, msg.y, msg.z

    raise RuntimeError("No LOCAL_POSITION_NED received")
\end{verbatim}

This is useful for reading the vehicle state, defining a home position, or generating waypoints relative to the current pose. 

\subsubsection{Sending a local position instruction}

A local position instruction is sent using \texttt{SET\_POSITION\_TARGET\_LOCAL\_NED}. The mask ignores velocity, acceleration, and yaw rate, meaning that only position and, optionally, yaw are commanded.

\begin{verbatim}
def send_local_position(master, x_m, y_m, z_down_m, yaw_rad=None):
    type_mask = (
        mavutil.mavlink.POSITION_TARGET_TYPEMASK_VX_IGNORE |
        mavutil.mavlink.POSITION_TARGET_TYPEMASK_VY_IGNORE |
        mavutil.mavlink.POSITION_TARGET_TYPEMASK_VZ_IGNORE |
        mavutil.mavlink.POSITION_TARGET_TYPEMASK_AX_IGNORE |
        mavutil.mavlink.POSITION_TARGET_TYPEMASK_AY_IGNORE |
        mavutil.mavlink.POSITION_TARGET_TYPEMASK_AZ_IGNORE |
        mavutil.mavlink.POSITION_TARGET_TYPEMASK_YAW_RATE_IGNORE
    )

    if yaw_rad is None:
        type_mask |= mavutil.mavlink.POSITION_TARGET_TYPEMASK_YAW_IGNORE
        yaw_rad = 0.0

    master.mav.set_position_target_local_ned_send(
        0,
        master.target_system,
        master.target_component,
        mavutil.mavlink.MAV_FRAME_LOCAL_NED,
        type_mask,
        x_m, y_m, z_down_m,
        0, 0, 0,
        0, 0, 0,
        yaw_rad,
        0
    )
\end{verbatim}

Since the frame is \texttt{LOCAL\_NED}, a positive altitude above the vehicle origin must be sent as a negative down value. For example, flying at \(1.5\) m above the local origin means sending:

\[
z_{\text{down}} = -1.5
\]

This sign conversion is also used directly in the supplied code through:
\begin{verbatim}
z_down = -z_up_m
\end{verbatim}

\subsubsection{Holding position by repeatedly sending the same command}

In guided mode, a single setpoint is usually not enough. The autopilot expects repeated updates, so the scripts send the same local position command at a fixed rate. In the provided code, this is done using a helper function that transmits setpoints at a configurable frequency, typically \(10\) Hz.

\begin{verbatim}
def hold_position(master, x_m, y_m, z_up_m, yaw_rad, duration_s, rate_hz=10.0):
    import time
    dt = 1.0 / rate_hz
    z_down = -z_up_m
    n = max(1, int(duration_s * rate_hz))

    for _ in range(n):
        send_local_position(master, x_m, y_m, z_down, yaw_rad)
        time.sleep(dt)
\end{verbatim}

A simple usage example is:

\begin{verbatim}
hold_position(master, x_m=1.0, y_m=0.5, z_up_m=1.5,
              yaw_rad=None, duration_s=5.0, rate_hz=10.0)
\end{verbatim}

This commands the vehicle to hold the point \((1.0, 0.5, -1.5)\) in \texttt{LOCAL\_NED} for five seconds.

\subsubsection{Landing}

Landing can be triggered by changing the flight mode:

\begin{verbatim}
set_mode(master, "LAND")
\end{verbatim}

\subsubsection{Minimal runnable example}

The following compact example combines the essential steps into one script:

\begin{verbatim}
from pymavlink import mavutil
import time

master = mavutil.mavlink_connection('udpin:127.0.0.1:14551')
master.wait_heartbeat()

def set_mode(mode_name):
    mode_map = master.mode_mapping()
    mode_id = mode_map[mode_name]
    master.mav.set_mode_send(
        master.target_system,
        mavutil.mavlink.MAV_MODE_FLAG_CUSTOM_MODE_ENABLED,
        mode_id
    )

def arm():
    master.mav.command_long_send(
        master.target_system,
        master.target_component,
        mavutil.mavlink.MAV_CMD_COMPONENT_ARM_DISARM,
        0, 1, 0, 0, 0, 0, 0, 0
    )

def takeoff(alt):
    master.mav.command_long_send(
        master.target_system,
        master.target_component,
        mavutil.mavlink.MAV_CMD_NAV_TAKEOFF,
        0, 0, 0, 0, 0, 0, 0, alt
    )

def send_local_position(x, y, z_down, yaw=None):
    type_mask = (
        mavutil.mavlink.POSITION_TARGET_TYPEMASK_VX_IGNORE |
        mavutil.mavlink.POSITION_TARGET_TYPEMASK_VY_IGNORE |
        mavutil.mavlink.POSITION_TARGET_TYPEMASK_VZ_IGNORE |
        mavutil.mavlink.POSITION_TARGET_TYPEMASK_AX_IGNORE |
        mavutil.mavlink.POSITION_TARGET_TYPEMASK_AY_IGNORE |
        mavutil.mavlink.POSITION_TARGET_TYPEMASK_AZ_IGNORE |
        mavutil.mavlink.POSITION_TARGET_TYPEMASK_YAW_RATE_IGNORE
    )

    if yaw is None:
        type_mask |= mavutil.mavlink.POSITION_TARGET_TYPEMASK_YAW_IGNORE
        yaw = 0.0

    master.mav.set_position_target_local_ned_send(
        0,
        master.target_system,
        master.target_component,
        mavutil.mavlink.MAV_FRAME_LOCAL_NED,
        type_mask,
        x, y, z_down,
        0, 0, 0,
        0, 0, 0,
        yaw,
        0
    )

set_mode("GUIDED")
time.sleep(2)
arm()
time.sleep(2)
takeoff(1.5)
time.sleep(8)

for _ in range(50):
    send_local_position(1.0, 0.0, -1.5, None)
    time.sleep(0.1)

set_mode("LAND")
\end{verbatim}


\subsubsection{Communication Logic and Port Usage}

In this setup, MAVLink communication is handled through two UDP endpoints with different roles.

The first endpoint is used by MAVProxy:
\begin{verbatim}
mavproxy.py --master :14550 --console --map
\end{verbatim}

Here, port \texttt{14550} is the main MAVLink input seen by MAVProxy. In other words, MAVProxy receives the telemetry stream arriving from the flight controller through the ESP-based Wi-Fi link on this port.

The Python control script does not connect directly to this same endpoint. Instead, it opens a second MAVLink connection:

\begin{verbatim}
udpin:127.0.0.1:14551
\end{verbatim}

This means the script listens locally on port \texttt{14551} for a MAVLink stream that is made available on the PC. In practice, MAVProxy acts as middleware: it receives the original telemetry stream and exposes a second local endpoint that the Python script can use.

Therefore, the port roles can be summarized as:

\begin{itemize}
    \item \texttt{14550}: MAVProxy master/input port
    \item \texttt{14551}: local MAVLink endpoint used by the Python script
\end{itemize}

This is why two different port numbers appear in the workflow. They are not arbitrary duplicates, but separate stages in the communication chain.

The full logic is:

\begin{enumerate}
    \item The flight controller is connected to the ESP module.
    \item The ESP module creates a Wi-Fi hotspot.
    \item The ground station computer connects to this hotspot.
    \item Telemetry arrives to MAVProxy on port \texttt{14550}.
    \item MAVProxy exposes a local MAVLink stream on port \texttt{14551}.
    \item The Python script connects to port \texttt{14551}, receives telemetry, and sends control commands through that connection.
\end{enumerate}

This arrangement is useful because MAVProxy can serve as an intermediate communication layer, while the Python code remains focused on vehicle control logic.
\subsubsection{Running the Code}

Before running the Python control script, the following conditions should be satisfied:

\begin{enumerate}
    \item The flight controller must be powered and connected to the ESP module.
    \item The ESP module must be active and broadcasting its Wi-Fi hotspot.
    \item The computer running the script must be connected to that Wi-Fi network.
    \item MAVLink telemetry must be available on the expected UDP port.
    \item The autopilot must support the selected guided control mode.
\end{enumerate}

Once the wireless telemetry link is available, the script can be started normally from the terminal. A minimal example is:

\begin{verbatim}
python3 your_script.py
\end{verbatim}

\subsection{General Remarks}

Using \texttt{LOCAL\_NED} control provides several advantages:

\begin{itemize}
    \item High precision in indoor or GPS-denied environments
    \item Simple trajectory definition in Cartesian space
    \item Compatibility with external localization systems (e.g., motion capture or SLAM)
\end{itemize}

However, care must be taken with:
\begin{itemize}
    \item Sign conventions (especially for the $z$ axis)
    \item Consistent reference frame definition
    \item Maintaining continuous command streaming
\end{itemize}

\subsection{Further Reading}

A detailed description of MAVLink commands in guided mode can be found in the official ArduPilot documentation:

\begin{center}
\texttt{https://ardupilot.org/dev/docs/copter-commands-in-guided-mode.html}
\end{center}